En una prueba sintética de secado de un recubrimiento, el soporte y la materia seca pesan juntos \(20\,\mathrm g\). El conjunto pesa \(32\,\mathrm g\) al entrar en el secador y \(28\,\mathrm g\) dos minutos después. Se mantienen constantes la temperatura del aire y su circulación; la única masa que puede salir es solvente. Se quiere retirar la pieza cuando pese \(24\,\mathrm g\).

Elija el estado y proponga una ley diferencial de secado que considere defendible, explicitando sus supuestos. Calíbrela y calcule cuándo retiraría la pieza. Decida si los datos hacen única esa previsión y justifique su respuesta mediante otra elección razonable de modelo. Proponga una nueva lectura que ayude a decidir entre sus propuestas e indique qué comprobaría con ella.
